\subsection{Comparison with an arbitrary actual uniformization}
\label{subsec:ftb-arbitrary-complex-uniformization}

Keep the actual complex principally polarized abelian variety \(P\),
its line bundle \(L\), its rigidification, and its dimension \(g\).
Assume the five algebraic laws retained in
Proposition~\ref{prop:ftb-full-fixed-fibre-theta-realization}.
That proposition gives one uniformization
\(u_0=(e_0,\Lambda_0,(f_{0,i}))\), its Appell--Humbert data
\((H_0,\chi_0)\), a perfect integral alternating form
\(E_0=\operatorname{Im}H_0|_{\Lambda_0}\), and the specified
bijection between analytic vectors and algebraic tangent vectors.
In this subsection \(u=(e,\Lambda,(f_i))\) is any
\(\mathrm{HolUnif}(P,L,g)\).
In particular \(e:\mathbb C^g\to P(\mathbb C)\) is an additive
surjection, its kernel is the discrete full lattice \(\Lambda\),
regular-function evaluations along \(e\) are holomorphic on
their actual inverse-image opens, and the \(f_i\) are the
nonvanishing analytic frames of the same algebraic line \(L\).
These conditions do not initially include a tangent isomorphism
for \(e\).

\begin{srlemma}[The source's linear lift]
\label{lem:ftb-holomorphic-linear-lift}
Let \(A/\mathbb C\) be an abelian variety, and let
\(e_0:\mathbb C^g\to A(\mathbb C)\) be a supplied
\(\mathrm{HolUnif}\) with its two tangent-bijection clauses.
If \(G:\mathbb C^m\to A(\mathbb C)\) is additive and
regular-evaluation-holomorphic, there is a unique complex
linear \(F:\mathbb C^m\to\mathbb C^g\) with
\(e_0F=G\).
\end{srlemma}
\begin{proof}
The proof uses separating regular functions and Baire's theorem;
it does not assume that \(G\) already has a holomorphic lift.
Choose an affine open \(U\) containing the identity of \(A\),
put \(S=U(\mathbb C)\), and use all \(f\in\Gamma(U,\mathcal O_U)\)
as separating complex-valued functions on \(S\).
They separate its actual complex points by the affine
point--ring-homomorphism identification.
Their inverse-image domains for \(e_0\) and \(G\) are open.
All their compositions are holomorphic.

Here is the weak openness assertion needed to lift.
For every neighborhood \(N\) of \(0\) in \(\mathbb C^g\),
\[
                 G(z)\in e_0(N)\qquad\text{for all sufficiently small }z.
                                                               \tag{AU1}
\]
Choose \(R>0\) with
\(\overline B(0,R)\subset N\cap e_0^{-1}(S)\), and put \(r=R/4\).
Let
\[
 T=\{z:\exists w\in\overline B(0,2r),\ e_0(w)=G(z)\}.
\]
Then \(T\subset G^{-1}(S)\), and
\(\overline T\cap G^{-1}(S)\subset T\).
Indeed, from \(z_j\to z\) and witnesses \(w_j\) in the indicated
compact ball, extract a convergent subsequence \(w_j\to w\).
Continuity of every separating regular-function evaluation
gives \(e_0(w)=G(z)\).

Choose a countable dense sequence \(c_j\) in \(\mathbb C^g\)
and set
\[
 B_j=\{z:\exists w\in B(c_j,r),\ e_0(w)=G(z)\}.
\]
Surjectivity of \(e_0\) gives
\(\bigcup_j\overline{B_j}=\mathbb C^m\).
Baire's theorem gives a \(j\) for which
\(\overline{B_j}\) has nonempty interior.
Choose \(z_0\in B_j\cap\operatorname{int}\overline{B_j}\)
and a witness \(w_0\in B(c_j,r)\).
If \(z\in B_j\), subtraction of the witnesses gives
\(z-z_0\in T\).
Thus, for \(z\) sufficiently near \(0\),
\(z+z_0\in\overline{B_j}\), hence \(z\in\overline T\).
Also \(G(z)\in S\) near \(0\).
The preceding relative-closedness assertion yields \(z\in T\),
proving (AU1).

Discreteness of \(\ker e_0\) gives a neighborhood \(W\) of \(0\)
on which \(e_0\) is injective: take \(W-W\) meeting
\(\ker e_0\) only at \(0\).
Let \(s\) be its inverse on \(e_0(W)\).
Assertion (AU1), applied to arbitrarily small neighborhoods
inside \(W\), shows that \(h=sG\) is defined near \(0\),
\(h(0)=0\), and \(h\) is continuous at \(0\).
Near \((0,0)\),
\[
                    h(x+y)-h(x)-h(y)\in\ker e_0.
\]
This expression tends to \(0\).
Discreteness makes it zero on a neighborhood of \((0,0)\).
The resulting locally additive germ extends uniquely to a
continuous real linear map \(F_{\mathbb R}\):
for fixed \(v\) use
\(2^j h(2^{-j}v)\) for sufficiently large \(j\).
Local additivity makes this eventually independent of \(j\);
continuity of the germ makes the resulting additive map continuous.
The identity \(e_0F_{\mathbb R}=G\) holds near \(0\), hence
everywhere by additivity and dyadic division.

It remains to prove complex linearity.
For any \(f\in\Gamma(U,\mathcal O_U)\), write
\(a_f(w)=f(e_0(w))\), \(b_f(z)=f(G(z))\).
The real chain rule in \(a_fF_{\mathbb R}=b_f\), together
with the complex differentiability of \(a_f,b_f\), gives
\[
       d a_f(0)\bigl(F_{\mathbb R}(cv)-cF_{\mathbb R}(v)\bigr)=0
                         \quad(c\in\mathbb C,\ v\in\mathbb C^m).
                                                               \tag{AU2}
\]
The differentials \(da_f(0)\) jointly have zero kernel.
To see this from the stated hypotheses, associate to a vector \(w\)
its algebraic dual-number tangent vector by the second
tangent-bijection clause for \(e_0\).
If every \(da_f(0)(w)\) vanishes, its dual-number homomorphism
on the affine chart equals that of the zero tangent vector.
The first tangent-bijection clause then gives \(w=0\).
Apply this to (AU2).
Thus \(F_{\mathbb R}\) is complex linear.
If two real linear lifts exist, their difference maps every
real line into the discrete subgroup \(\ker e_0\); continuity
on that line makes the difference zero.
This proves the stated uniqueness as well.
\end{proof}

\begin{srlemma}[A nowhere-zero Appell--Humbert function]
\label{lem:ftb-trivial-appell-humbert-data}
Let \(\Lambda\subset\mathbb C^g\) be a discrete full real lattice.
Let \(B\) be Hermitian, let \(|c_\lambda|=1\), and let \(r\)
be nowhere-zero entire, with
\[
 r(z+\lambda)=c_\lambda
       \exp\!\left(\pi B(z,\lambda)+\frac{\pi}{2}B(\lambda,\lambda)\right)r(z).
                                                               \tag{AU3}
\]
Then \(B=0\), all \(c_\lambda=1\), and \(r\) is constant.
\end{srlemma}
\begin{proof}
The function
\[
       \Phi(z)=|r(z)|\exp\!\left(-\frac{\pi}{2}B(z,z)\right)
\]
is continuous, positive and \(\Lambda\)-periodic.
A compact fundamental parallelepiped bounds both \(\Phi\)
and \(\Phi^{-1}\).
Consequently, for constants \(M,M'\),
\[
 |r(z)|\le M e^{\pi B(z,z)/2},\qquad
 |r(z)|^{-1}\le M'e^{-\pi B(z,z)/2}.                         \tag{AU4}
\]
If \(B(v,v)<0\), the entire function \(t\mapsto r(tv)\)
is bounded and tends to \(0\) along real \(t\to+\infty\).
Liouville's theorem makes it constant, equal to \(0\), a contradiction.
If \(B(v,v)>0\), apply the same argument to \(t\mapsto r(tv)^{-1}\).
Thus \(B(v,v)=0\) for every \(v\).
Expanding at \(x+y\) and \(x+iy\) gives \(B(x,y)=0\), so \(B=0\).
The first bound in (AU4) now makes \(r\) bounded, hence constant.
Substitute \(z=0\) in (AU3) and cancel \(r(0)\ne0\)
to obtain \(c_\lambda=1\).
\end{proof}

\begin{srproposition}[The supplied frame has the same tangent and integral form]
\label{prop:ftb-any-uniformization-tangent-perfect}
For the arbitrary actual \(u\) above there is a unique complex
linear isomorphism \(F:\mathbb C^g\overset{\sim}\longrightarrow\mathbb C^g\)
with \(e_0F=e\); it restricts to a lattice isomorphism
\(\phi:\Lambda\overset{\sim}\longrightarrow\Lambda_0\).
Every algebraic tangent vector has a unique analytic vector
for \(u\), and every analytic vector gives a unique algebraic
tangent vector.
If the actual frames of \(u\) have Appell--Humbert data
\((H,\chi)\), then
\[
       H_0(Fz,Fw)=H(z,w),\qquad
       E(a,b)=E_0(\phi a,\phi b)=\operatorname{Im}H(a,b).       \tag{AU5}
\]
In particular \(E:\Lambda\to\operatorname{Hom}_{\mathbb Z}(\Lambda,\mathbb Z)\)
is bijective.
The tangent correspondence respects the ordered linear
combinations used in the key datum.
\end{srproposition}
\begin{proof}
Apply Lemma~\ref{lem:ftb-holomorphic-linear-lift} to \(G=e\).
The kernel identities give
\[
                         Fz\in\Lambda_0\iff z\in\Lambda.       \tag{AU6}
\]
If \(Fz=0\), every \(tz\), \(t\in\mathbb C\), belongs to
\(\Lambda\).
The continuous map \(t\mapsto tz\) into this discrete space
is constant, so \(z=0\).
Finite equal complex dimensions make \(F\) invertible.
Equation (AU6) then gives the asserted lattice isomorphism.

On the inverse-image of a line chart \(U_i\), form
\[
                             r_i(z)=\frac{f_{0,i}(Fz)}{f_i(z)}.
\]
The same algebraic transition functions occur above and below;
their evaluated factors cancel on overlaps.
The \(r_i\) therefore glue to a nowhere-zero entire \(r\).
Their actual frame automorphy gives (AU3) with
\[
 B(z,w)=H(z,w)-H_0(Fz,Fw),\qquad
 c_\lambda=\chi(\lambda)/\chi_0(F\lambda).
\]
Here the semicharacters have norm one.
Lemma~\ref{lem:ftb-trivial-appell-humbert-data} gives \(B=0\).
This proves the first identity in (AU5).
The pullback \(\phi^\vee E_0\phi\) is integral and bijective,
and its values are the second expression in (AU5).

For a regular function on any actual open containing the
identity, the identity \(e=e_0F\) and the complex chain rule give
\[
                   u.\mathrm{IsDiff}(w,v)
                      \iff u_0.\mathrm{IsDiff}(Fw,v).          \tag{AU7}
\]
The inverse chain rule uses \(F^{-1}\).
The two tangent-bijection clauses for \(u_0\) therefore
give both clauses for \(u\).
Finally, if \(v=\sum_j a_jv_j\) in the algebraic Lie space and
\(v_j\) corresponds to \(w_j\), then each regular-function
dual-number derivative of \(v\) is the sum of those of the
\(a_jv_j\).
The same is true of the derivative of \(e\) at
\(\sum_j a_jw_j\).
Equality of the dual-number homomorphisms on an affine
identity chart, followed by localization, proves that
\(v\) corresponds to \(\sum_j a_jw_j\).
Thus (AU7) is compatible with the actual ordered Lie basis
and its row-matrix changes, not just with abstract dimensions.
\end{proof}

For later comparisons the same argument has the exact
power version used in the source.
If \(u,u'\) uniformize actual lines \(L,L'\), an algebraic map
\(f:P\to P'\) is lifted by \(F\), and
\[
                         (f^*L')^{\otimes m}\simeq L^{\otimes n},
\]
evaluate this line isomorphism in its actual two chart frames.
Its evaluations \(w_{ij}(e(z))\) satisfy the evaluated
transition identities, so
\[
        r_{ij}(z)=w_{ij}(e(z))
                      \frac{f'_j(Fz)^m}{f_i(z)^n}
\]
glue to a nowhere-zero entire function.
Apply Lemma~\ref{lem:ftb-trivial-appell-humbert-data} to
\(B=nH-mF^*H'\) and
\(c_\lambda=\chi(\lambda)^n/\chi'(F\lambda)^m\).
It gives, with exactly these tensor-power exponents,
\[
                              mH'(Fz,Fw)=nH(z,w).             \tag{AU8}
\]
In particular no analytic line comparison has been substituted
for the given algebraic line isomorphism.

\paragraph{Frozen source and scope.}
All sources below are at \srcid{alpoge-source-7cbf7a0c}.
The linear-lift mechanism is
\srcid{src/Development/AbelianVarietyTheory/Proofs/CubicalLineBundles/ExistsUniqueLinearLiftOfHolomorphicAbelianVariet.lean}
(\srcid{5f25708bfa7aa5609008d34be077a14b4d21af81}),
in particular \srcid{tendsto_map_nhds_of_separating} and
\srcid{existsUnique_linear_lift_of_separating}.
Its real additive-germ extension and discrete-difference uniqueness
are
\srcid{src/Development/AbelianVarietyTheory/Proofs/HomogeneousIdeals/ExistsUniqueRealLinearOfEventuallyAddModDiscrete.lean}
(\srcid{0e6a11a999b466dd1c459c4d2089b0c999350d44}).
The Gaussian argument is
\srcid{src/Development/AbelianVarietyTheory/Proofs/CubicalLineBundles/AppellHumbertTrivialOfNowhereVanishing.lean}
(\srcid{63ce58372c5f22dfd3c30c464bd9290d251da6bf}).
The evaluated line-isomorphism construction, its transition
cancellation, and (AU8) are
\srcid{HermFormPullEqOfLbIsoPow.lean}
(\srcid{8da182d30841698bb02f4773f5df334e7b3396c2})
in that same directory.
The actual map/tangent chain rule is
\srcid{ExistsLinearOfIsMonHom.lean}
(\srcid{d3df5da8674f36fd9e33ebc871ed5b08454e4daf}).
The two assembly forms are
\srcid{HolomorphicUniformizationTangentUnimodularRiemann.lean}
(\srcid{9e0d2e677fce33d4c5c91b778c4fe70b845c3cad})
and
\srcid{src/Development/AbelianVarietyTheory/Proofs/ThetaLawProofs/UnimodularAndTangentOfIsAH.lean}
(\srcid{a4df7948258a8441bcf1d2ffcd6db4a65fb44846}).

This establishes the actual \srcid{UnimodularAndTangentOfIsAH}
interface from the five retained algebraic laws and the fixed-fibre
construction.  It supplies the arbitrary-frame integral form and
the differential normalization needed for the Gaussian
archimedean norm calculation.
It does not construct an algebraic variety from an arbitrary
period matrix, nor a family of uniformizations; those are separate
subsequent producers.
